\documentclass{article}
\usepackage{amsmath, amssymb, amsthm}

\setlength{\parindent}{0pt}

\newtheorem{claim}{Claim}
\newtheorem*{corollary}{Corollary}

\begin{document}

In a certain card game, each card has a point value.

\begin{itemize}
    \item Numbered cards in the range $2$ to $9$ are worth five points each.
    \item The card numbered $10$ and the face cards (jack, queen, king) are worth ten points each.
    \item Aces are worth fifteen points each.
\end{itemize}

\begin{claim}
    Suppose that you thoroughly shuffle a $52$-card deck. The expected total point value of the three 
    cards on the top of the deck after the shuffle is $\frac{285}{13}$.
\end{claim}

\begin{proof}
    For each $i \in \{1, 2, 3\}$, define the random variable 
    \[
    X_i =
    \begin{cases}
        5 & \text{if card } i \text{ is worth five points} \\
        10 & \text{if card } i \text{ is worth ten points} \\
        15 & \text{if card } i \text{ is worth fifteen points}
    \end{cases}
    \]
    Let $X$ denote the total point value of the first three cards. Then
    \[
    X = X_1 + X_2 + X_3
    \]
    By linearity of expectation,
    \[
    \mathrm{E}[X] = \mathrm{E}[X_1] + \mathrm{E}[X_2] + \mathrm{E}[X_3]
    \]
    Since the random variables $X_1$, $X_2$, and $X_3$ have the same probability distribution, 
    \[
    \mathrm{E}[X] = 3\mathrm{E}[X_1]
    \]
    so it remains only to compute $\mathrm{E}[X_1]$. By the definition of expectation, 
    \[
    \mathrm{E}[X_1] = 5 \cdot \Pr\{X_1 = 5\} + 10 \cdot \Pr\{X_1 = 10\} + 15 \cdot \Pr\{X_1 = 15\}
    \]
    The relevant probabilities are
    \[
    \begin{aligned}
        \Pr\{X_1 = 5\} &= \frac{32}{52} = \frac{8}{13} \\[4pt]
        \Pr\{X_1 = 10\} &= \frac{16}{52} = \frac{4}{13} \\[4pt]
        \Pr\{X_1 = 15\} &= \frac{4}{52} = \frac{1}{13}
    \end{aligned}
    \]
    Substituting the above values into the above equation gives 
    \[
    \mathrm{E}[X_1] = 5 \cdot \frac{8}{13} + 10 \cdot \frac{4}{13} + 15 \cdot \frac{1}{13} = \frac{95}{13}
    \]
    Thus 
    \[
    \mathrm{E}[X] = \frac{95}{13} + \frac{95}{13} + \frac{95}{13} = \frac{285}{13}
    \]

    \medskip

    Hence the expected total point value of the first three cards is $\frac{285}{13}$ points.
\end{proof}

\begin{claim}
    Suppose that you throw out all the red cards and shuffle the remaining $26$-card, all-black deck. 
    The expected total point value of the top three cards is $\frac{285}{13}$.
\end{claim}

\begin{proof}
    For each $i \in \{1, 2, 3\}$, define the random variable 
    \[
    X_i =
    \begin{cases}
        5 & \text{if card } i \text{ is worth five points} \\
        10 & \text{if card } i \text{ is worth ten points} \\
        15 & \text{if card } i \text{ is worth fifteen points}
    \end{cases}
    \]
    Let $X$ denote the total point value of the first three cards. Then
    \[
    X = X_1 + X_2 + X_3
    \]
    By linearity of expectation,
    \[
    \mathrm{E}[X] = \mathrm{E}[X_1] + \mathrm{E}[X_2] + \mathrm{E}[X_3]
    \]
    Since the random variables $X_1$, $X_2$, and $X_3$ have the same probability distribution, 
    \[
    \mathrm{E}[X] = 3\mathrm{E}[X_1]
    \]
    so it remains only to compute $\mathrm{E}[X_1]$. By the definition of expectation, 
    \[
    \mathrm{E}[X_1] = 5 \cdot \Pr\{X_1 = 5\} + 10 \cdot \Pr\{X_1 = 10\} + 15 \cdot \Pr\{X_1 = 15\}
    \]
    The relevant probabilities are
    \[
    \begin{aligned}
        \Pr\{X_1 = 5\} &= \frac{16}{26} = \frac{8}{13} \\[4pt]
        \Pr\{X_1 = 10\} &= \frac{8}{26} = \frac{4}{13} \\[4pt]
        \Pr\{X_1 = 15\} &= \frac{2}{26} = \frac{1}{13}
    \end{aligned}
    \]
    Substituting the above values into the above equation gives 
    \[
    \mathrm{E}[X_1] = 5 \cdot \frac{8}{13} + 10 \cdot \frac{4}{13} + 15 \cdot \frac{1}{13} = \frac{95}{13}
    \]
    Thus 
    \[
    \mathrm{E}[X] = \frac{95}{13} + \frac{95}{13} + \frac{95}{13} = \frac{285}{13}
    \]

    \medskip

    Hence the expected total point value of the first three cards is $\frac{285}{13}$ points.
\end{proof}

\begin{corollary}
    Removing all cards of either color does not change the expected total point value of the top 
    three cards, i.e., $\frac{285}{13}$.
\end{corollary}

\begin{proof}
    Claim $2$ establishes the case where the red cards are removed, leaving an all-black deck.

    \medskip

    The case where the black cards are removed follows by symmetry. Each of the $13$ ranks has exactly $2$ 
    red and $2$ black cards, so interchanging the roles of red and black in the proof of Claim $2$ gives the 
    corresponding result for the all-red deck.

    \medskip

    Hence the expected total point value of the top three cards is $\frac{285}{13}$ regardless.
\end{proof}

\end{document}
